\documentclass[11pt,a4paper]{article}
% Spolecny styl pro ucitelske pripravy MAT (Work/MAT)
% Verze 1.1, 2026-10-23
\NeedsTeXFormat{LaTeX2e}
\ProvidesPackage{priprava}[2026/10/23 v1.1 Sablona ucitelskych priprav MAT]

% Zakladni typografie a cestina
\RequirePackage[T1]{fontenc}
\RequirePackage[utf8]{inputenc}
\RequirePackage[czech]{babel}
\RequirePackage{lmodern}
\RequirePackage{microtype}
\RequirePackage{amsmath,amssymb}
\RequirePackage{xcolor}
\RequirePackage{enumitem}
\RequirePackage[most]{tcolorbox}
\RequirePackage{tabularx,booktabs,array}
\RequirePackage{fancyhdr}
\RequirePackage{needspace}
\RequirePackage{geometry}
\RequirePackage[hidelinks]{hyperref}

\geometry{a4paper,top=17mm,bottom=18mm,left=23mm,right=23mm}
\definecolor{notegray}{gray}{0.39}
\setlength{\parindent}{0pt}
\setlength{\parskip}{2.5pt plus 1pt minus 0.5pt}
\setlist[itemize]{topsep=2pt,itemsep=1pt,parsep=0pt,leftmargin=1.45em}
\setlist[enumerate]{topsep=2pt,itemsep=2pt,parsep=0pt,leftmargin=1.65em}
\setcounter{secnumdepth}{0}
\renewcommand{\arraystretch}{1.12}

% Zahlavi: vlevo rocnik | globalni tema, vpravo cislo: nazev hodiny
% Zapati zustava zcela prazdne.
\pagestyle{fancy}
\fancyhf{}
\renewcommand{\headrulewidth}{0.35pt}
\renewcommand{\footrulewidth}{0pt}
\setlength{\headheight}{14pt}
\newcommand{\nastavPripravu}[4]{%
  \fancyhead[L]{\small #1\enspace\textbar\enspace #2}%
  \fancyhead[R]{\small #3: #4}%
}

% Minimalisticke cernobile boxy. Kazda definice zvlast.
\tcbset{commonlesson/.style={enhanced,breakable=false,colback=white,colframe=black,
  colbacktitle=white,coltitle=black,boxrule=0.45pt,
  arc=0.8mm,outer arc=0.8mm,left=3mm,right=3mm,top=1.6mm,bottom=1.6mm,
  toptitle=1mm,bottomtitle=0.7mm,fonttitle=\bfseries,
  attach boxed title to top left={xshift=2.5mm,yshift=-1.15mm},
  boxed title style={colback=white,colframe=white,boxrule=0pt,left=1mm,right=1mm}}}
\newtcolorbox{zadaniBox}[1][]{commonlesson,title={Zadání},#1}
\newtcolorbox{prikladBox}[1][]{commonlesson,title={Příklad},#1}
\newtcolorbox{definiceBox}[1][]{commonlesson,title={Definice},#1}

% Spolecny box pro zadani a jeho reseni, bez mezery mezi castmi.
% Pouziti:
%   \begin{prikladReseniBox}
%     Zadani ulohy.
%     \tcblower
%     \begin{align*} ... \end{align*}
%   \end{prikladReseniBox}
% Spodni cast ma jemny sedy podklad a delici linku; lze v ni sazet
% vice radku matematiky. Obe casti maji jen jeden spolecny ramecek.
\tcbset{solutionlesson/.style={
  commonlesson,
  enhanced jigsaw,
  breakable,
  colbacklower=black!3,
  segmentation style={solid,draw=black!25,line width=0.35pt},
  fontlower=\small\color{notegray},
  before lower={\noindent\textbf{Řešení:}\enspace},
  leftlower=3mm,rightlower=3mm
}}
\newtcolorbox{prikladReseniBox}[1][]{solutionlesson,title={Příklad},#1}
\newtcolorbox{zadaniReseniBox}[1][]{solutionlesson,title={Zadání},#1}

% Pomucky pro autora pripravy: sedive poznamky; \reseni pro starsi pripravy
\newcommand{\poznamka}[1]{\par\begingroup\small\color{notegray}\itshape
  \textnormal{\textbf{Poznámka pro učitele.}}\enspace #1\par\endgroup}
\newcommand{\reseni}[1]{\par\begingroup\small\color{notegray}
  \textbf{Očekávané řešení:}\enspace #1\par\endgroup}
\newcommand{\blok}[1]{\needspace{5\baselineskip}\subsection*{#1}}
\newcommand{\vyrok}[1]{\uv{#1}}

\nastavPripravu{1. ročník}{Algebraické výrazy}{48}{Dělení mnohočlenu mnohočlenem}
\begin{document}

% Zdroje: vlastni ucebnice Algebraicke vyrazy, Deleni mnohoclenu mnohočlenem;
% Realisticky.cz: 010805 Dělení mnohočlenů (metodická inspirace).
\poznamka{Jádro 45 minut: analogie s písemným dělením čísel 5 min, algoritmus na tabuli 12 min, společný příklad 10 min, samostatná práce 13 min, kontrola 5 min. Dělení se zbytkem je podle tempa závěr nebo rezerva.}
\blok{Písemné dělení čísel a mnohočlenů má stejnou logiku}
Při dělení odhadujeme nejvyšší řád, zpětně násobíme a odečítáme. U mnohočlenů místo nejvyšších cifer sledujeme \textbf{členy s nejvyšší mocninou}. Oba mnohočleny nejdřív seřadíme podle klesajících exponentů.
\begin{zadaniBox}[title={Podíl a zbytek při dělení mnohočlenů}]
Pro mnohočleny $A(x)$ a $B(x)\ne0$ existují jednoznačné mnohočleny $Q(x)$ (\textbf{podíl}) a $R(x)$ (\textbf{zbytek}), pro které
\[A(x)=B(x)Q(x)+R(x),\]
přičemž $R=0$ nebo $\deg R<\deg B$. Pro podíl výrazů $A(x)/B(x)$ navíc požadujeme $B(x)\ne0$.
\end{zadaniBox}
\blok{Algoritmus: vyděl, vynásob, odečti, pokračuj}
\begin{prikladReseniBox}[breakable=false,title={Příklad 1: dělení beze zbytku}]
Vyděl $(x^2+5x+6):(x+2)$.
\tcblower
Nejvyšší členy: $x^2:x=x$. Zpětné násobení a odečtení:
\[(x^2+5x+6)-(x^2+2x)=3x+6.\]
Dále $3x:x=3$ a $(3x+6)-3(x+2)=0$. Proto $Q(x)=x+3$, $R(x)=0$. Zkouška: $(x+2)(x+3)=x^2+5x+6$. Podíl výrazů platí pro $x\ne-2$.
\end{prikladReseniBox}
\begin{prikladReseniBox}[breakable=false,title={Příklad 2: záporné koeficienty}]
Vyděl $(2x^3-x^2-13x+5):(2x+5)$.
\tcblower
\begin{align*}
2x^3:(2x)&=x^2,\\
(2x^3-x^2-13x+5)-x^2(2x+5)&=-6x^2-13x+5,\\
(-6x^2):(2x)&=-3x,\\
(-6x^2-13x+5)-(-3x)(2x+5)&=2x+5,\\
(2x):(2x)&=1.
\end{align*}
Po posledním odečtení dostaneme $0$. Podíl $Q(x)=x^2-3x+1$, zbytek $R=0$. Podíl výrazů je definován pro $x\ne-\frac52$.
\end{prikladReseniBox}
\poznamka{Dbej na znak mínus před \emph{celým} zpětně násobeným mnohočlenem. Než žák začne další krok, musí správně sečíst zbývající členy. Zkoušku dělej vždy součinem dělitele a podílu, případně přičtením zbytku.}
\blok{Samostatná práce — základ}
\begin{zadaniReseniBox}[breakable=false,title={Samostatná práce: urči podíl a zbytek}]
\begin{enumerate}[label=\alph*)]
\item $(x^2+7x+10):(x+5)$;
\item $(x^2-4):(x-2)$;
\item $(2x^2+7x+3):(2x+1)$;
\item $(x^3-1):(x-1)$.
\end{enumerate}
\tcblower
\begin{enumerate}[label=\alph*)]
\item $Q=x+2$, $R=0$; zkouška $(x+5)(x+2)$;
\item $Q=x+2$, $R=0$; zkouška $(x-2)(x+2)$;
\item $Q=x+3$, $R=0$; zkouška $(2x+1)(x+3)$;
\item $Q=x^2+x+1$, $R=0$; zkouška $(x-1)(x^2+x+1)$.
\end{enumerate}
\end{zadaniReseniBox}
\blok{Zbytek, který nezmizí}
\begin{prikladReseniBox}[breakable=false,title={Příklad 3: dělení se zbytkem}]
Vyděl $(x^2+2x+3):(x+1)$.
\tcblower
$x^2:x=x$; po odečtení $x(x+1)$ zůstává $x+3$. Další člen podílu je $1$; po odečtení $(x+1)$ zůstává $2$. Proto $Q=x+1$, $R=2$ a
\[x^2+2x+3=(x+1)(x+1)+2.\]
Podíl výrazů lze zapsat $x+1+\frac2{x+1}$ pro $x\ne-1$.
\end{prikladReseniBox}
\blok{Volitelná rezerva}
\begin{zadaniReseniBox}[breakable=false,title={Rozšíření: chybějící mocnina}]
Vyděl $(x^3+2x+1):(x+1)$. Nezapomeň na nulový koeficient u $x^2$.
\tcblower
$Q=x^2-x+3$, $R=-2$; zkouška $(x+1)(x^2-x+3)-2=x^3+2x+1$.
\end{zadaniReseniBox}
\blok{Domácí procvičení}
\begin{zadaniReseniBox}[breakable=false,title={Domácí úkol: podíl, zbytek a zkouška}]
\begin{enumerate}[label=\alph*)]
\item $(x^2+3x+2):(x+1)$;
\item $(x^2+1):(x-1)$;
\item $(2x^2+5x+2):(2x+1)$;
\item $(x^3+1):(x+1)$;
\item $(x^3+2x+1):(x+1)$;
\item $(3x^2+7x+2):(x+2)$.
\end{enumerate}
\tcblower
\begin{enumerate}[label=\alph*)]
\item $Q=x+2$, $R=0$;
\item $Q=x+1$, $R=2$;
\item $Q=x+2$, $R=0$;
\item $Q=x^2-x+1$, $R=0$;
\item $Q=x^2-x+3$, $R=-2$;
\item $Q=3x+1$, $R=0$.
\end{enumerate}
\end{zadaniReseniBox}
\end{document}
